% Part: first-order-logic
% Chapter: sequent-calculus
% Section: propositional-rules

\documentclass[../../../include/open-logic-section]{subfiles}

\begin{document}

\olfileid[hi]{mvl}{seq}{prl}

\olsection{चयनित तर्कशास्त्रों के प्रतिज्ञप्तिक नियम}

किसी $n$-पक्षीय सीक्वेंट कलन में किसी संयोजक के अनुमान-नियम केवल उस संयोजक
के अभिलाक्षणिक सत्य-फलन पर निर्भर करते हैं। अतः यदि विभिन्न तर्कशास्त्रों
में कोई संयोजक एक ही सत्य-फलन से परिभाषित हो, तो उस संयोजक के ये $n$-पक्षीय
सीक्वेंट नियम उन तर्कशास्त्रों में समान होते हैं।

\subsection{$\lnot$ के नियम}

$\lnot$ के निम्नलिखित नियम लुकासिएविच और क्लिनी तर्कशास्त्रों तथा उनके
रूपांतरों पर लागू होते हैं।

\begin{defish}
  \begin{center}
\AxiomC{$ \Gamma \nSequent \Pi \nSequent \Delta, !A $}
\RightLabel{\iR{\lnot}{\False}}
\UnaryInfC{$ \lnot !A, \Gamma \nSequent \Pi \nSequent \Delta$}
\DisplayProof
\\[2ex]
\AxiomC{$ \Gamma \nSequent !A, \Pi \nSequent \Delta $}
\RightLabel{\iR{\lnot}{\Undef}}
\UnaryInfC{$\Gamma \nSequent \lnot !A, \Pi \nSequent \Delta$}
\DisplayProof
\\[2ex]
\AxiomC{$!A, \Gamma \nSequent \Pi \nSequent \Delta$}
\RightLabel{\iR{\lnot}{\True}}
\UnaryInfC{$\Gamma \nSequent \Pi \nSequent \Delta,  \lnot !A$}
\DisplayProof
  \end{center}
\end{defish}

$\lnot$ के निम्नलिखित नियम गोडेल तर्कशास्त्र पर लागू होते हैं।

\begin{defish}
\AxiomC{$ \Gamma \nSequent !A, \Pi \nSequent \Delta, !A $}
\RightLabel{\iR{\lnot}{\False}[\LogGod]}
\UnaryInfC{$ \lnot !A, \Gamma \nSequent \Pi \nSequent \Delta$}
\DisplayProof
\hfill
\AxiomC{$!A, \Gamma \nSequent \Pi \nSequent \Delta$}
\RightLabel{\iR{\lnot}{\True}[\LogGod]}
\UnaryInfC{$\Gamma \nSequent \Pi \nSequent \Delta,  \lnot !A$}
\DisplayProof
\hfill
\end{defish}

(गोडेल तर्कशास्त्र में $\lnot !A$ कभी $\Undef$ मूल्य नहीं ले सकता, इसलिए
मध्य स्थान के लिए कोई नियम नहीं है।)

\subsection{$\land$ के नियम}

ये लुकासिएविच, प्रबल क्लिनी और गोडेल तर्कशास्त्र में $\land$ के नियम हैं।

\begin{defish}
\begin{center}
  \AxiomC{$!A, !B, \Gamma \nSequent \Pi \nSequent \Delta$}
  \RightLabel{$\iR{\land}{\False}$}
  \UnaryInfC{$!A \land !B, \Gamma \nSequent \Pi \nSequent \Delta$}
  \DisplayProof
\\[2ex]
  \AxiomC{$\Gamma \nSequent !A, \Pi \nSequent !A, \Delta$}
  \AxiomC{$\Gamma \nSequent !B, \Pi \nSequent !B, \Delta$}
  \AxiomC{$\Gamma \nSequent !A, !B, \Pi \nSequent \Delta$}
  \RightLabel{$\iR\land\Undef$}
  \TrinaryInfC{$\Gamma \nSequent !A \land !B, \Pi \nSequent \Delta$}
  \DisplayProof
  \\[2ex]  
  \AxiomC{$\Gamma \nSequent \Pi \nSequent \Delta, !A$}
  \AxiomC{$\Gamma \nSequent \Pi \nSequent \Delta, !B$}
  \RightLabel{$\iR\land\True$}
  \BinaryInfC{$\Gamma \nSequent \Pi \nSequent \Delta, !A \land !B$}
  \DisplayProof
\end{center}
\end{defish}

\subsection{$\lor$ के नियम}

ये लुकासिएविच, प्रबल क्लिनी और गोडेल तर्कशास्त्र में $\lor$ के नियम हैं।

\begin{defish}
\begin{center}
 \AxiomC{$!A, \Gamma \nSequent \Pi \nSequent \Delta$}
  \AxiomC{$!B, \Gamma \nSequent \Pi \nSequent \Delta$}
  \RightLabel{$\iR\lor\False$}
  \BinaryInfC{$!A \lor !B, \Gamma \nSequent \Pi \nSequent \Delta$}
  \DisplayProof
\\[2ex]
  \AxiomC{$!A, \Gamma \nSequent !A, \Pi \nSequent \Delta$}
  \AxiomC{$!B, \Gamma \nSequent !B, \Pi \nSequent \Delta$}
  \AxiomC{$\Gamma \nSequent !A, !B, \Pi \nSequent \Delta$}
  \RightLabel{$\iR\lor\Undef$}
  \TrinaryInfC{$\Gamma \nSequent !A \lor !B, \Pi \nSequent \Delta$}
  \DisplayProof
  \\[2ex]  
  \AxiomC{$\Gamma \nSequent \Pi \nSequent \Delta, !A, !B$}
  \RightLabel{$\iR\lor\True$}
  \UnaryInfC{$\Gamma \nSequent \Pi \nSequent \Delta, !A \lor !B$}
  \DisplayProof
\end{center}
\end{defish}

\subsection{$\lif$ के नियम}

ये लुकासिएविच तर्कशास्त्र में $\lif$ के नियम हैं।

\begin{defish}
\begin{center}
  \AxiomC{$\Gamma \nSequent \Pi \nSequent \Delta, !A$}
  \AxiomC{$!B, \Gamma \nSequent \Pi \nSequent \Delta$}
  \RightLabel{$\iR\lif\False[\LogLuk[3]]$}
  \BinaryInfC{$!A \lif !B, \Gamma \nSequent \Pi \nSequent \Delta$}
  \DisplayProof
  \\[2ex]
  \AxiomC{$\Gamma \nSequent !A, !B, \Pi \nSequent \Delta$}
  \AxiomC{$!B, \Gamma \nSequent \Pi \nSequent \Delta, !A$}
  \RightLabel{$\iR\lif\Undef[\LogLuk[3]]$}
  \BinaryInfC{$\Gamma \nSequent !A \lif !B, \Pi \nSequent \Delta$}
  \DisplayProof
\\[2ex]
  \AxiomC{$!A, \Gamma \nSequent !B, \Pi \nSequent \Delta, !B$}
  \AxiomC{$!A, \Gamma \nSequent !A, \Pi \nSequent \Delta, !B$}
  \RightLabel{$\iR{\lif}{\True}[\LogLuk[3]]$}
  \BinaryInfC{$\Gamma \nSequent \Pi \nSequent \Delta, !A \lif !B$}
  \DisplayProof
\end{center}
\end{defish}

ये प्रबल क्लिनी तर्कशास्त्र में $\lif$ के नियम हैं।

\begin{defish}
\begin{center}
  \AxiomC{$\Gamma \nSequent \Pi \nSequent \Delta, !A$}
  \AxiomC{$!B, \Gamma \nSequent \Pi \nSequent \Delta$}
  \RightLabel{$\iR\lif\False[\LogKs]$}
  \BinaryInfC{$!A \lif !B, \Gamma \nSequent \Pi \nSequent \Delta$}
  \DisplayProof
  \\[2ex]
  \AxiomC{$!B, \Gamma \nSequent !B, \Pi \nSequent \Delta$}
  \AxiomC{$\Gamma \nSequent !A, !B, \Pi \nSequent \Delta$}
  \AxiomC{$\Gamma \nSequent !A, \Pi \nSequent \Delta, !A$}
  \RightLabel{$\iR\lif\Undef[\LogKs]$}
  \TrinaryInfC{$\Gamma \nSequent !A \lif !B, \Pi \nSequent \Delta$}
  \DisplayProof
\\[2ex]
  \AxiomC{$!A, \Gamma \nSequent \Pi \nSequent \Delta, !B$}
  \RightLabel{$\iR{\lif}{\True}[\LogKs]$}
  \UnaryInfC{$\Gamma \nSequent \Pi \nSequent \Delta, !A \lif !B$}
  \DisplayProof
\end{center}
\end{defish}

ये गोडेल तर्कशास्त्र में $\lif$ के नियम हैं।

\begin{defish}
\begin{center}
  \AxiomC{$\Gamma \nSequent !A, \Pi \nSequent \Delta, !A$}
  \AxiomC{$!B, \Gamma \nSequent \Pi \nSequent \Delta$}
  \RightLabel{$\iR\lif\False[\LogGod[3]]$}
  \BinaryInfC{$!A \lif !B, \Gamma \nSequent \Pi \nSequent \Delta$}
  \DisplayProof
  \\[2ex]
  \AxiomC{$\Gamma \nSequent !B, \Pi \nSequent \Delta$}
  \AxiomC{$\Gamma \nSequent \Pi \nSequent \Delta, !A$}
  \RightLabel{$\iR\lif\Undef[\LogGod[3]]$}
  \BinaryInfC{$\Gamma \nSequent !A \lif !B, \Pi \nSequent \Delta$}
  \DisplayProof
\\[2ex]
  \AxiomC{$!A, \Gamma \nSequent !B, \Pi \nSequent \Delta, !B$}
  \AxiomC{$!A, \Gamma \nSequent !A, \Pi \nSequent \Delta, !B$}
  \RightLabel{$\iR{\lif}{\True}[\LogGod[3]]$}
  \BinaryInfC{$\Gamma \nSequent \Pi \nSequent \Delta, !A \lif !B$}
  \DisplayProof
\end{center}
\end{defish}


\begin{sidewaysfigure}
  \begin{center}  
  \AxiomC{$A \nSequent A \nSequent A$}
  \RightLabel{\iR \Weakening \True}
  \UnaryInfC{$A \nSequent A \nSequent B, A$}
  \RightLabel{\iR \Weakening \Undef}
  \UnaryInfC{$A \nSequent B, A \nSequent B, A$}
  \RightLabel{\iR \Weakening \Undef}
  \UnaryInfC{$A \nSequent A, B, A \nSequent B, A$}
  \AxiomC{$A \nSequent A \nSequent A$}
  \RightLabel{\iR \Weakening \True}
  \UnaryInfC{$A \nSequent A \nSequent A, A$}
  \RightLabel{\iR \Weakening \True}
  \UnaryInfC{$A \nSequent A \nSequent B, A, A$}
  \RightLabel{\iR \Weakening \False}
  \UnaryInfC{$B, A \nSequent A \nSequent B, A, A$}
  \RightLabel{$\iR\lif\Undef$}
  \BinaryInfC{$A \nSequent A \lif B, A \nSequent B, A$}
  \AxiomC{$B \nSequent B \nSequent B$}
  \RightLabel{\iR \Weakening \Undef}
  \UnaryInfC{$B \nSequent A, B \nSequent B$}
  \RightLabel{\iR \Exchange \Undef}
  \UnaryInfC{$B \nSequent B, A \nSequent B$}
  \RightLabel{\iR \Weakening \Undef}
  \UnaryInfC{$B \nSequent A, B, A \nSequent B$}
  \RightLabel{\iR \Weakening \False}
  \UnaryInfC{$A, B \nSequent A, B, A \nSequent B$}
  \RightLabel{\iR \Exchange \False}
  \UnaryInfC{$B, A \nSequent A, B, A \nSequent B$}
  \AxiomC{$A \nSequent A \nSequent A$}
  \RightLabel{\iR \Weakening \True}
  \UnaryInfC{$A \nSequent A \nSequent B, A$}
  \RightLabel{\iR \Weakening \False}
  \UnaryInfC{$B, A \nSequent A \nSequent B, A$}
  \RightLabel{\iR \Weakening \False}
  \UnaryInfC{$B, B, A \nSequent A \nSequent B, A$}
  \RightLabel{$\iR\lif\Undef$}
  \BinaryInfC{$B, A \nSequent A \lif B, A \nSequent B$}
  \RightLabel{$\iR\lif\False$}
  \BinaryInfC{$A \lif B, A \nSequent A \lif B, A \nSequent B$}
  \DisplayProof
  \end{center}
  \caption{$\LogLuk[3]$ में !!{derivation} का उदाहरण}
\end{sidewaysfigure}

\end{document}
